% docs/texte/05_theory/imaging.tex -- theory chapter "Imaging" (Interferosim)
% Version: 0.1.0 (2026-10-12)
% Notation as defined in texte/05_theory/notation.tex.
% ================================================================

\index{dirty image}\index{dirty beam}\index{weighting}
\section{Dirty image and dirty beam}

The measured visibilities are $\mathcal S(u,v)\,\mathcal V(u,v)$. Their
inverse Fourier transform, with imaging weights $W(u,v)$, is the
\emph{dirty image}
\begin{equation}
  I_D(l,m) = \iint W\,\mathcal S\,\mathcal V\;
             e^{+2\pi i\,(ul+vm)}\,\mathrm{d}u\,\mathrm{d}v
           = (A\,I) * B_D ,
\end{equation}
the true (beam-weighted) sky convolved with the \emph{dirty beam}
\begin{equation}
  B_D(l,m) = \iint W\,\mathcal S\;e^{+2\pi i\,(ul+vm)}\,\mathrm{d}u\,\mathrm{d}v .
\end{equation}
$B_D$ is the point spread function (PSF) of the array. Normalising
$B_D(0,0)=1$ gives $I_D$ in Jy\,beam$^{-1}$: a point source shows its
flux density as peak value.

\section{Gridding}

To use the FFT, the visibilities are placed on a regular grid of
$N\times N$ cells. For an image of $N$ pixels of size $\Delta l$ the
cell size in the $(u,v)$ plane is $\Delta u = 1/(N\,\Delta l)$.
Interferosim currently uses nearest-cell gridding (each sample is added
to the nearest cell); convolutional gridding is a possible extension.

\section{Weighting}

The weights $W$ trade sensitivity against beam shape:
\begin{itemize}
  \item \textbf{Natural weighting}: every sample has the same weight
        (or $1/\sigma^2$). Best signal-to-noise ratio, but densely
        sampled regions dominate and the beam has higher sidelobes.
  \item \textbf{Uniform weighting}: every occupied cell has the same
        weight, $W_k = 1/N_k$ for $N_k$ samples in cell $k$. Better
        beam, more noise.
  \item \textbf{Robust (Briggs) weighting} \cite{briggs1995}
        interpolates continuously between both.
\end{itemize}
The Wiener weighting of chapter~\ref{ch:deconvolution} is another
interpolation of the same kind.
